\begin{solapartat}
\punts{0,2} La condició d'escapament exigeix que la $E_m = 0$

\punts{0,5} $E_m = E_C + E_P = 0 \Rightarrow \dfrac{1}{2}mv^2_{esc} + \left(-G\dfrac{Mm}{R}\right) = 0$

\punts{0,3} $\dfrac{1}{2}mv^2_{esc} = G\dfrac{Mm}{R} \Rightarrow v_{esc} = \sqrt{\dfrac{2GM}{R}}$
\end{solapartat}

\begin{solapartat}
Denominem, $v_{\mathit{inicial}} = v_0 = \dfrac{1}{2}v_{esc} = \dfrac{1}{2}\sqrt{\dfrac{2GM}{R_L}}$

\punts{0,1} Abans de tornar a caure, la $v_{\mathit{final}} = 0$

\punts{0,1} $E_{m,\mathit{inicial}} = E_{m,\mathit{final}}$

\punts{0,2} $\dfrac{1}{2}mv_0^2 + \left(-G\dfrac{Mm}{R_L}\right) = 0 + \left(-G\dfrac{Mm}{R}\right)$; on $R = R_L + h$ i $h$ l'altura respecte a la superfície lunar, a la qual arribarà l'objecte abans de tornar a caure.

\punts{0,4}
\begin{align*}
&m\left(\frac{1}{2}\,\frac{1}{4}\,\frac{2GM}{R_L} - \frac{GM}{R_L}\right) = -G\frac{M}{R} \Rightarrow\\
&\Rightarrow \frac{1}{4}\,\frac{1}{R_L} - \frac{1}{R_L} = -\frac{1}{R} \Rightarrow\\
&\Rightarrow -\frac{3}{4}\,\frac{1}{R_L} = -\frac{1}{R} \Rightarrow R = \frac{4}{3}R_L
\end{align*}

\punts{0,2} $R = \dfrac{4}{3}1737\times 10^{3} = 2,316\times 10^{6}\ m \Rightarrow h = R - R_L = 579\ km$
\end{solapartat}
